\usepackage{xcolor}


% =============================================================================
% Базовые цвета
% =============================================================================

\definecolor{xblue}{HTML}{3182CE}
\definecolor{xgray}{HTML}{718096}
\definecolor{xgreen}{HTML}{38A169}
%\definecolor{xindigo}{HTML}{4C51BF}
\definecolor{xmoss}{HTML}{5D715E}
\definecolor{xolive}{HTML}{5F7A1E}
\definecolor{xorange}{HTML}{DD6B20}
\definecolor{xpink}{HTML}{D53F8C}
%\definecolor{xpurple}{HTML}{805AD5}
%\definecolor{xred}{HTML}{BD4242}
\definecolor{xrust}{HTML}{9C4221} 
\definecolor{xsepia}{HTML}{8B5E3C}
\definecolor{xstone}{HTML}{999999}


\definecolor{xpurple}{HTML}{6B46C1}   % чуть темнее и менее «электрический»
\definecolor{xred}{HTML}{B03030}      % холоднее, лучше держит насыщенность
\definecolor{xindigo}{HTML}{3C40A0}


% =============================================================================
% Цвета определений, утверждений, лемм, теорем, и т.п.
% =============================================================================

\colorlet{deficolor}{xblue}
\colorlet{statcolor}{xgreen!85!black}
\colorlet{lemmcolor}{xgray}
\colorlet{theocolor}{xpurple}
\colorlet{conccolor}{xolive!65!black}
\colorlet{examcolor}{xsepia!80!white}
\colorlet{notecolor}{xpink!80}
\colorlet{algocolor}{xrust!40!black}
\colorlet{codecolor}{xindigo!30!black}


% =============================================================================
% Цвета списков
% =============================================================================

\colorlet{listL1}{xred!70!black}
\colorlet{listL2}{xred!55!black}
\colorlet{listL3}{xred!40!black}


% =============================================================================
% Цвета todo-шек
% =============================================================================

\colorlet{todobackground}{xorange!20}     % Светло-оранжевый фон
\colorlet{todoborder}{xorange}            % Оранжевая рамка


% =============================================================================
% Специальные цвета
% =============================================================================

\colorlet{xcovere}{xmoss}                 % Цвет на обложке
\colorlet{xmatrix}{xmoss}                 % Базовый «матричный» зелёный
\colorlet{xmatrixDark}{xmoss!75!black}    % Тёмный оттенок для текста
\colorlet{xmatrixGlow}{xmoss!40}          % Полупрозрачный для «свечения»